\Article{8}
The following objections can be made to this proof.

1. The rule by which Euler concludes that all $2m-2$ unknowns $\alpha,\beta$, etc., $\lambda,\mu$, etc. can be determined rationally from $2m-1$ equations is by no means general, but very often admits exceptions. If, for example, in Article 3 one regards one of the unknowns as known and tries to express the others rationally in terms of it and the given coefficients, one readily finds that this is impossible, and that no unknown quantity can be determined except by an equation of degree $m-1$. Although here one can immediately see a priori that this must necessarily happen, one might nevertheless reasonably doubt whether, in the present case too, matters stand in the same way for some values of $m$, so that the unknowns $\alpha,\beta$, etc., $\lambda,\mu$, etc. cannot be determined from $u,B,C$, etc. except by an equation of degree perhaps higher than $2m$. For the case where the equation \mbox{$X=0$} is of the fourth degree, Euler derives rational values of the coefficients in terms of $u$ and the given coefficients; but that this can also be done for all higher equations certainly needed a fuller explanation. --- It seems worthwhile, moreover, to investigate more deeply and in the greatest generality those formulas that express $\alpha,\beta$, etc. rationally in terms of $u,B,C$, etc. I shall undertake a fuller treatment of this and other matters belonging to the theory of elimination (a subject by no means exhausted) on another occasion.

2. Even if it has been proved that, whatever the degree of the equation \mbox{$X=0$}, formulas can always be found that express $\alpha,\beta$, etc., $\lambda,\mu$, etc. rationally in terms of $u,B,C$, etc., it is nevertheless certain that for some particular values of the coefficients $B,C$, etc. these formulas can become \textit{indeterminate}, so that not only is it impossible to determine these unknowns rationally from $u,B,C$, etc., but in some cases there are actually no real values of $\alpha,\beta$, etc., $\lambda,\mu$, etc. corresponding to a real value of $u$. To confirm this, for brevity I refer the reader to Euler's paper itself, where the equation of the fourth degree is discussed more fully on p. 236. Everyone will immediately see that the formulas for the coefficients $\alpha,\beta$ become indeterminate\SourceLeaf{019} if $C=0$ and the value $0$ is taken for $u$, and that their values not only cannot be assigned without extracting roots, but are not even real if $BB-4D$ is a negative quantity. Although in this case one can readily see that $u$ has other real values to which real values of $\alpha,\beta$ correspond, one might nevertheless fear that unraveling this difficulty (which Euler did not touch on at all) could require much greater effort in higher equations. Certainly this matter must by no means be passed over in silence in an exact proof.

3. The distinguished Euler tacitly supposes that the equation \mbox{$X=0$} has $2m$ roots, and puts their sum $=0$ because the second term is absent in $X$. I have already stated above, in Article 3, what I think of this freedom (which all authors on this subject use). The proposition that the sum of all the roots of an equation equals its first coefficient with its sign changed seems applicable only to equations that have roots. Since this very proof is meant to establish that the equation \mbox{$X=0$} actually has roots, it seems inadmissible to suppose their existence. No doubt those who have not yet seen through the fallacy in this reasoning will reply that \textit{what is proved here is not that the equation \mbox{$X=0$} can be satisfied} (for that is what the expression that it has roots means), \textit{but only that it can be satisfied by values of $x$ contained in the form $a+b\sqrt{-1}$; the former is supposed as an axiom}. But since no forms of quantities other than the real and the imaginary $a+b\sqrt{-1}$ can be conceived, it is not sufficiently clear how what is to be proved differs from what is supposed as an axiom. Indeed, even if it were possible to devise still other forms of quantities, say the form $F,F',F''$, etc., it should nevertheless not be admitted without proof that an equation can be satisfied by some value of $x$ that is either real, or contained in the form $a+b\sqrt{-1}$, or in the form $F$, or in $F'$, etc. That axiom can therefore have no other meaning than this: any equation can be satisfied \textit{either} by a real value of the unknown, \textit{or} by an imaginary value contained in the form $a+b\sqrt{-1}$, \textit{or} perhaps by a value contained in some other form hitherto unknown, \textit{or} by a value contained in no form at all. But how quantities of this kind, of which you cannot even form an idea --- a true shadow of a shadow --- can be added or multiplied is certainly not understood with the clarity always required in mathematics\HistoricalNote{4}.

\SourceLeaf{020}
I do not, however, wish these objections to cast any doubt on the conclusions Euler drew from his supposition. On the contrary, I am certain that they can be established by a method neither difficult nor very different from \textsc{Euler's}, in such a way that not even the slightest doubt should remain for anyone. I criticize only the \textit{form}, which, although it can be very useful in \textit{discovering} new truths, nevertheless seems quite unacceptable in \textit{proving} them before the public.

4. The distinguished Euler offered nothing at all to prove the assertion that the product $pqr$, etc. can be determined \textit{rationally} from the coefficients in $X$. Everything he says about this for equations of the \textit{fourth degree} is as follows (where $a,b,c,d$ are the roots of the proposed equation $x^4+Bxx+Cx+D=0$):

\lq No doubt it will be objected that I have supposed here that the quantity $pqr$ was a real quantity and that its square $ppqqrr$ was positive; this was still doubtful, since, the roots $a,b,c,d$ being imaginary, it could well happen that the square of the quantity $pqr$ formed from them was negative. To this I reply that this case can never occur; for however imaginary the roots $a,b,c,d$ may be, we nevertheless know that there must be $a+b+c+d=0$: $ab+ac+ad+bc+bd+cd=B$; $abc+abd+acd+bcd=-C$\HistoricalNote{5}; $abcd=D$, these quantities $B,C,D$ being real. But since $p=a+b$, $q=a+c$, $r=a+d$, their product $pqr=(a+b)(a+c)(a+d)$ can be determined, \textit{as we know}, from the quantities $B,C,D$, and will consequently be real, just as we have seen that in fact $pqr=-C$ and $ppqqrr=CC$. It will readily be recognized in the same way that this same circumstance must occur in higher equations, and that no objections can be made to me on this ground.\rq{}
Euler nowhere added the condition that the product $pqr$, etc. can be determined \textit{rationally} in terms of $B,C$, etc.; nevertheless, he seems always to have understood it, since without it the proof can have no force. Now it is indeed true for equations of the fourth degree that expanding the product $(a+b)(a+c)(a+d)$ gives $aa(a+b+c+d)+abc+abd+acd+bcd=-C$; yet it does not seem sufficiently clear how in all higher equations the product can be determined rationally\SourceLeaf{021} from the coefficients. The celebrated \textsc{de Foncenex}, who first observed this (\textit{Miscell. phil. math. soc. Taurin.}, Vol. I, p. 117), rightly maintains that without a rigorous proof of this proposition the method loses all its force; he confesses, however, that this proof seems rather difficult to him, and describes the path he tried in vain\HistoricalNote{6}. Nevertheless, the matter can be settled without much difficulty by the following method (whose substance I can only indicate here). Although in equations of the fourth degree it is not sufficiently clear that the product $(a+b)(a+c)(a+d)$ can be determined in terms of the coefficients $B,C,D$, one can readily see that the same product is also $=(b+a)(b+c)(b+d)$, as well as $=(c+a)(c+b)(c+d)$, and finally also $=(d+a)(d+b)(d+c)$. The product $pqr$ will therefore be one quarter of the sum $(a+b)(a+c)(a+d)+(b+a)(b+c)(b+d)+(c+a)(c+b)(c+d)+(d+a)(d+b)(d+c)$, which, when expanded, can readily be foreseen a priori to be a polynomial in the roots $a,b,c,d$ in which they all enter in the same way. Such functions can always be expressed rationally in terms of the coefficients of the equation whose roots are $a,b,c,d$. --- The same thing is also evident if the product $pqr$ is put into the form
\[
\tfrac12(a+b-c-d)\times\tfrac12(a+c-b-d)\times\tfrac12(a+d-b-c)
\]
for one can readily foresee that this product, when expanded, will involve all of $a,b,c,d$ in the same way. Those familiar with the subject will also readily gather from this how it must be applied to higher equations. --- A complete account of the proof, which brevity does not allow me to include here, together with a fuller investigation of functions involving several variables in the same way, I reserve for another occasion.

I observe, moreover, that besides these four objections there are still some other points in Euler's proof that could be criticized; I pass over them in silence, however, lest I seem too severe a critic, especially since the foregoing seems sufficiently to show that the proof, at least in the form Euler proposed it, can by no means be regarded as complete.

After this proof Euler showed another way of reducing the theorem for equations whose degree is not a power of two to the solution of equations of that kind. But since this method teaches nothing for equations whose
\SourceLeaf{022}
degree is a power of two, and is moreover open to all the preceding objections (except the fourth) just as the first general proof is, there is no need to discuss it more fully here.
